% Deutsche Parallelfassung zu state/part_intro.tex

\chapter*{Zustand}
\phantomsection
\addcontentsline{toc}{chapter}{Zustand}

Das Gerade--Übergangsbogen--Kreisbogen-Alignment aus dem vorangehenden Teil
bleibt der behandelte Engineering Object. Seine Identität, geordnete
Konstruktion, der intrinsische Parameter \(s\), das Krümmungsgesetz
\(\kappa(s)\), die Berlinish-Übergangsbogenpartition, Provenienz und
Anwendbarkeit sind bereits geklärt. Zustand beginnt deshalb nicht mit einem
anonymen Koordinatentupel, sondern mit einer Frage zu diesem identifizierten
Objekt:

\begin{quote}
Welcher Eisenbahnzustand folgt an einem gewählten \(s\) aus dem Alignment, und
welche zusätzlichen Gesetze und Realisierungsdaten werden dafür benötigt?
\end{quote}

Die Krümmung bestimmt zunächst die Richtungsänderung und mit einer Anfangspose
eine ebene Trajektorie. Höhengeometrie und Überhöhung ergänzen eigenständige
Längsgesetze. Eine angegebene Frame-Konvention verbindet diese Eingaben zu
einer räumlichen Eisenbahninterpretation. Erst mit Geschwindigkeit und, falls
erforderlich, einem Fahrzeugmodell folgen Beschleunigung und weitere
betriebliche Größen. Die zweckbezogene Bewertung liegt nochmals dahinter: Eine
berechnete Folge ist keine Engineering Decision.

\begin{figure}[htbp]
\centering
\begin{tikzpicture}[
  node distance=5mm and 5mm,
  box/.style={draw,rounded corners,align=center,minimum height=8mm,text width=25mm,font=\small},
  arrow/.style={->,thick}
]
\node[box] (id) {identifiziertes\\Alignment};
\node[box,right=of id] (law) {\(s,\ \kappa(s)\)\\intrinsisches Gesetz};
\node[box,right=of law] (p2) {Anfangspose\\+\ pose2};
\node[box,below=of p2] (frame) {bewegtes Frame\\+ Höhe + Überhöhung};
\node[box,left=of frame] (p3) {qualifizierter\\pose3-Zustand};
\node[box,left=of p3] (dyn) {Tempo + Modell\\abgeleitete Antwort};
\node[box,below=of dyn] (eval) {Zweck + Regel\\Bewertung};
\draw[arrow] (id)--(law);
\draw[arrow] (law)--(p2);
\draw[arrow] (p2)--(frame);
\draw[arrow] (frame)--(p3);
\draw[arrow] (p3)--(dyn);
\draw[arrow] (dyn)--(eval);
\node[align=center,font=\scriptsize,below=2mm of frame] {Realisierung, Konvention, Provenienz, Gültigkeit};
\end{tikzpicture}

\caption{Abhängigkeit vom identifizierten Alignment bis zu einer qualifizierten Engineering-Folge.}
\label{fig:railway-state-dependency}
\end{figure}

Die Pfeile in \cref{fig:railway-state-dependency} bezeichnen Abhängigkeit,
nicht Identität oder Autorität. Insbesondere wird eine realisierte Pose nicht
zur Alignment-Identität, eine Beobachtung nicht zur physischen Realität und
ein Bewertungsergebnis nicht zur Entscheidung über die Annahme eines
Vorschlags.

\section*{Ein enger Erklärungsvertrag}

Für diesen Teil kann ein \emph{Eisenbahnzustandsdatensatz} als folgende
erklärende Struktur der Thesis gelesen werden:
\[
\mathcal S_{\mathrm{rail}}(s)=
\bigl(
\mathrm{AlignmentRef},s,\mathrm{pose},F,\kappa,
\mathrm{vertical},\mathrm{cant},\mathrm{realization},
\mathrm{provenance},t_{\mathrm{valid}},\mathrm{derived}
\bigr).
\]
Geschwindigkeit gehört nur bei einer geschwindigkeitsqualifizierten Verwendung
dazu. Beobachtungsprovenienz gehört nur zu einer Schätzung des beobachteten
Zustands. Dieses Tupel ist weder ein neuer kanonischer Kernel-Begriff noch ein
universelles Speicherschema; es hält lediglich Abhängigkeiten und
Qualifikationen für den folgenden Argumentgang sichtbar.

\begin{table}[htbp]
\centering
\small
\begin{tabular}{p{0.18\textwidth}p{0.28\textwidth}p{0.43\textwidth}}
\toprule
\textbf{Art} & \textbf{Laufendes Beispiel} & \textbf{Was dadurch nicht feststeht}\\
\midrule
Intrinsisch & \(s\), Elementfolge, \(\kappa(s)\) & Koordinaten, Überhöhung, physischer Zustand, Annahme\\
Realisiert & lokale Pose, Höhe, Frame, dargestellte Koordinaten & Alignment-Identität oder tatsächlicher Zustand der gebauten Anlage\\
Beobachtet & zeitgestempelte gemessene Schienenlagen mit Methode und Unsicherheit & definierende Geometrie oder unqualifizierte Wahrheit\\
Abgeleitet & \(v^2\kappa\), effektive Querwirkung, Vergleichsresiduum & Autorität, Zulässigkeit oder Entscheidung\\
\bottomrule
\end{tabular}
\caption{Intrinsische, realisierte, beobachtete und abgeleitete Größen bleiben getrennt.}
\label{tab:railway-state-kinds}
\end{table}

Die Kapitel folgen nun einer einzigen Abhängigkeitskette:
\begin{align*}
\text{Alignment-Identität}&\rightarrow s\rightarrow\kappa(s)
\rightarrow\mathrm{pose2}\rightarrow\text{bewegtes Frame}\\
&\rightarrow\text{Höhe und Überhöhung}
\rightarrow\mathrm{pose3}\\
&\rightarrow\text{geschwindigkeitsabhängige Größen}
\rightarrow\text{zweckqualifizierte Bewertung}.
\end{align*}

\begin{principle}[Eisenbahnzustandsprinzip]
\label{princ:railway-state}
Eisenbahnzustand wird aus einem identifizierten Alignment unter angegebenen
zusätzlichen Gesetzen, Realisierungsdaten, Konventionen und
Verwendungsqualifikationen abgeleitet. Keine resultierende Pose,
Koordinatendarstellung, Beobachtung oder Bewertung besitzt die Identität oder
Autorität des Alignments.
\end{principle}

\begin{summary}
Teil V ist die physische Entfaltung des bereits identifizierten Alignments.
Jede Zustandsgröße tritt erst ein, wenn eine Engineering-Folge sie benötigt.
\end{summary}
